\chapter{التكاملات المنحنية والتكاملات المتعددة}\label{ch:b2:multint}

يوسّع هذا الفصل التكامل من الفترات إلى المنحنيات وإلى مناطق المستوي
والفضاء. فالتكاملات المنحنية تكامل \emph{\omterm{def:b2:multint:lineint}{صورة تفاضلية}} $P\,\dd x + Q\,\dd y$ على
طول قوس موجه؛ والتكاملات المزدوجة والثلاثية تكامل دوالا على مناطق
ذات بعدين وثلاثة أبعاد. وتلتقي النظريتان في \emph{مبرهنة غرين--ريمان}،
وهي المبرهنة الأساسية للتحليل ذات البعدين، أما الأداة الحسابية
الرئيسية في كل ما يلي فهي \emph{صيغة تغيير المتغيرات}، التي يكون
عامل تشوهها هو القيمة المطلقة \omterm{def:b2:linalg:det}{للمحدد} الجاكوبي.

\section{التكاملات المنحنية}

\begin{definition}[الصورة التفاضلية؛ التكامل المنحني]\label{def:b2:multint:lineint}
ليكن $U \subseteq \R^2$ مفتوحا. \emph{الصورة التفاضلية}\index{صورة تفاضلية} من
الدرجة $1$ والصنف $\mathcal{C}^0$ على $U$ هي عبارة $\omega = P\,\dd x + Q\,\dd y$ حيث $P, Q \colon U \to \R$ متصلتان
--- وبصورة رسمية، تطبيق \omterm{def:b2:metric:continuity}{متصل} من $U$ نحو ثنوي $\R^2$، أي $\omega(M) = P(M)\,e_1^* + Q(M)\,e_2^*$.
ومن أجل قوس $\gamma \colon [a, b] \to U$ من الصنف $\mathcal{C}^1$، $\gamma(t) =
(x(t), y(t))$، يكون \emph{التكامل
المنحني}\index{تكامل منحني} للصورة $\omega$ على طول $\gamma$ هو
\[
\int_\gamma \omega
= \int_a^b \Bigl(P(\gamma(t))\,x'(t)
+ Q(\gamma(t))\,y'(t)\Bigr)\,\dd t .
\]
وتمتد التعاريف حرفيا إلى $\R^3$ (الصور $P\,\dd x +
Q\,\dd y + R\,\dd z$) وإلى الأقواس من
الصنف $\mathcal{C}^1$ بالقطع (بالجمع على القطع).
\end{definition}

\begin{proposition}[الثبات والتوجيه]\label{prop:b2:multint:lineinv}
لا يتغير \omterm{def:b2:multint:lineint}{التكامل المنحني} عند تغيير وسيط متزايد من الصنف $\mathcal{C}^1$،
ويتغير إشارةً عند تغيير وسيط متناقص. ومنه فهو لا يتعلق إلا \omterm{def:b2:curves:reparam}{بالقوس
الهندسي} \emph{الموجه}.
\end{proposition}

\begin{proof}
إذا كان $\theta \colon [c, d] \to [a, b]$ تغيير وسيط و$\tilde\gamma = \gamma \circ \theta$، فإن قاعدة السلسلة وتغيير المتغير
بمتغير واحد $t = \theta(u)$ يعطيان
\[
\int_{\tilde\gamma}\omega
= \int_c^d \bigl(P(\gamma(\theta(u)))\,x'(\theta(u))
+ Q(\gamma(\theta(u)))\,y'(\theta(u))\bigr)\,\theta'(u)\,\dd u
= \pm\int_a^b \bigl(Px' + Qy'\bigr)(t)\,\dd t ,
\]
بالإشارة $+$ إذا كانت $\theta$ متزايدة ($\theta(c) = a$) والإشارة $-$ إذا
كانت متناقصة (فيتبادل الحدان).
\end{proof}

\begin{example}[عمل قوة؛ الدوران]\label{ex:b2:multint:work}
إذا كان $F = (P, Q)$ حقل قوة، فإن $\int_\gamma P\dd x + Q\dd y =
\int_a^b \langle F(\gamma(t)), \gamma'(t)\rangle\,\dd t$ هو \emph{عمل} $F$ على طول
$\gamma$. ومن أجل $\omega = -y\,\dd x +
x\,\dd y$ على دائرة الوحدة المقطوعة عكس عقارب الساعة
$\gamma(t) =
(\cos t, \sin t)$:
\[
\int_\gamma \omega
= \int_0^{2\pi}\bigl((-\sin t)(-\sin t)
+ \cos t\cos t\bigr)\,\dd t = 2\pi ,
\]
أي ضعف \omterm{def:b2:multint:domain}{المساحة} المحصورة --- وهو أول تلميح إلى مبرهنة غرين--ريمان.
\end{example}

\begin{example}[تكامل واحد، وتوسيمان، وفخ إشارة واحد]\label{ex:b2:multint:semicircle}
احسب $\int_\gamma x\,\dd y$ على طول نصف دائرة الوحدة العلوي من $(1, 0)$ إلى $(-1, 0)$. مع
$\gamma(t) =
(\cos t, \sin t)$، $t \in \intcc0\pi$:
\[
\int_0^\pi\cos t\cdot\cos t\,\dd t = \frac\pi2 .
\]
وبتوسيم التمثيل البياني $x \mapsto (x, \sqrt{1 -
x^2})$، $x$ من $1$ إلى $-1$ (انتبه
للاتجاه!):
\[
\int_1^{-1}x\cdot\frac{-x}{\sqrt{1 - x^2}}\,\dd x
= \int_{-1}^{1}\frac{x^2}{\sqrt{1 - x^2}}\,\dd x
= \frac\pi2
\]
(ويختزله $x = \sin u$ إلى تكامل واليس). القيمة نفسها، كما تضمن \cref{prop:b2:multint:lineinv} ---
لكن فقط لأن المسارين يسيران من $(1,0)$ إلى $(-1,0)$؛ فعكس السير يقلب
الإشارة. وإغلاق المسار على طول المحور $x$ (حيث $\dd y = 0$) لا يضيف
شيئا، والمجموع $\frac\pi2$ هو \omterm{def:b2:multint:domain}{مساحة} نصف القرص: وهي أول حالة من صيغ \omterm{def:b2:multint:domain}{مساحة}
الحد في مبرهنة غرين--ريمان أدناه.
\end{example}

\begin{definition}[الصور التامة والصور المغلقة]\label{def:b2:multint:exact}
تكون الصورة $\omega = P\,\dd x + Q\,\dd y$ من الصنف $\mathcal{C}^0$ \emph{تامة}\index{صورة تامة} على
$U$ إذا وُجدت $f \in \mathcal{C}^1(U)$ (وهي \emph{كمون}\index{كمون}) تحقق $\omega = \dd f$،
أي $P = f_x$ و$Q
= f_y$. وتكون صورة من الصنف $\mathcal{C}^1$ \emph{مغلقة}\index{صورة مغلقة}
إذا كان $P_y = Q_x$ على $U$.
\end{definition}

\begin{theorem}[المبرهنة الأساسية للتكاملات المنحنية]\label{thm:b2:multint:ftc}
إذا كانت $\omega = \dd f$ تامة وكان $\gamma$ قوسا من الصنف $\mathcal{C}^1$ بالقطع في
$U$ من $A$ إلى $B$، فإن
\[
\int_\gamma \omega = f(B) - f(A) .
\]
وبوجه خاص يكون تكامل \omterm{def:b2:multint:exact}{صورة تامة} على أي قوس مغلق معدوما، وتكون كل \omterm{def:b2:multint:exact}{صورة
تامة} من الصنف $\mathcal{C}^1$ مغلقة.
\end{theorem}

\begin{proof}
$\frac{\dd}{\dd t}f(\gamma(t)) = f_x(\gamma(t))x'(t) +
f_y(\gamma(t))y'(t)$ بقاعدة السلسلة (\cref{ch:b2:diffcalc})، ومنه فإن الدالة تحت التكامل في \cref{def:b2:multint:lineint}
هي مشتقة $t \mapsto
f(\gamma(t))$، وتعطي المبرهنة الأساسية للتحليل النتيجة على كل قطعة؛
وتتلاشى القيم الوسطى تلسكوبيا. أما انغلاق الصور التامة من الصنف
$\mathcal{C}^1$ فهو مبرهنة شوارتز: $P_y =
f_{xy} = f_{yx} = Q_x$.
\end{proof}

\begin{example}[إعادة بناء كمون]\label{ex:b2:multint:potential}
لتكن $\omega = y\,\eu^{xy}\,\dd x + (x\,\eu^{xy} + 2y)\,\dd y$ على $\R^2$. وهي مغلقة: فالمشتقتان المتقاطعتان تساويان
معا $\eu^{xy}(1 + xy)$. ولإيجاد \omterm{def:b2:multint:exact}{كمون}، نكامل $P$ بدلالة $x$ عند تثبيت $y$:
\[
f(x, y) = \int y\,\eu^{xy}\,\dd x = \eu^{xy} + c(y),
\]
ثم نضبط $c$ بمطابقة $f_y$: فيعطي $x\,\eu^{xy} + c'(y) =
x\,\eu^{xy} + 2y$ المقدار $c(y) = y^2$. ومنه
$f(x,y) = \eu^{xy} +
y^2$، ومن أجل أي قوس من الصنف $\mathcal C^1$ بالقطع من $(0,0)$ إلى $(1,1)$،
\[
\int_\gamma\omega = f(1,1) - f(0,0) = (\eu + 1) - 1 = \eu ,
\]
بصورة مستقلة عن المسار --- والوصفة ذات الخطوتين (كامل بدلالة $x$،
وصحّح بدلالة $y$) هي العكس العملي للمبرهنة \cref{thm:b2:multint:ftc} على المناطق
التي تكون فيها الصور المغلقة تامة.
\end{example}

\begin{example}[المغلقة لا تعني التامة]\label{ex:b2:multint:angleform}
على $U = \R^2 \setminus \{0\}$، تكون \emph{صورة الزاوية}
\[
\omega = \frac{-y\,\dd x + x\,\dd y}{x^2 + y^2}
\]
مغلقة (بحساب مباشر: فكل من $P_y$ و$Q_x$ يساوي $\frac{y^2 - x^2}{(x^2+y^2)^2}$)، لكن
تكاملها على طول دائرة الوحدة هو $2\pi \neq 0$ (بالحساب نفسه في \cref{ex:b2:multint:work}،
مقسوما على $1$): ومنه فإن $\omega$ ليست تامة على $U$. ومحليا،
$\omega = \dd\theta$ من أجل تعيين $\theta$ للزاوية القطبية؛ والفشل شامل --- فلا
يمكن تعريف الزاوية تعريفا متصلا حول الثقب. وعلى المناطق الخالية من
الثقوب يختفي المرض: فعلى مفتوح \emph{نجمي الشكل}، تكون كل \omterm{def:b2:multint:exact}{صورة مغلقة}
من الصنف $\mathcal{C}^1$ تامة (مبرهنة بوانكاريه، \cref{exo:b2:multint:8}).
\end{example}

\section{التكاملات المزدوجة}

نسلّم بنظرية تكامل ريمان بمتغير واحد (مجلد السنة الأولى، و\cref{ch:b2:integration})
ونرسم صيغتها ذات المتغيرين. فللدالة $f$ المتصلة على مستطيل
$R = [a, b] \times [c, d]$ تكامل مزدوج $\iint_R f$، يُعرَّف بمجاميع ريمان على شبكات تماما كما
في حالة متغير واحد، ويُحسب بالتكرار:

\begin{theorem}[فوبيني على مستطيل]\label{thm:b2:multint:fubini}
من أجل $f$ متصلة على $R = [a,b] \times [c,d]$،
\[
\iint_R f
= \int_a^b \Bigl(\int_c^d f(x, y)\,\dd y\Bigr)\dd x
= \int_c^d \Bigl(\int_a^b f(x, y)\,\dd x\Bigr)\dd y .
\]
\end{theorem}

\begin{proof}
نضع $F(x) = \int_c^d f(x, y)\,\dd y$. ويجعل الاتصال المنتظم للدالة $f$ على المتراص
$R$ الدالة $F$ متصلة (بالتقدير المهيمن: $\abs{F(x) - F(x')} \leq (d - c)\sup_y\abs{f(x,y) - f(x',y)}$).
ولنجزّئ الآن $[a,b]$ و$[c,d]$ إلى $n$ جزءا متساويا، فنحصل على شبكة
من الخلايا $R_{ij}$ \omterm{def:b2:multint:domain}{مساحة} كل منها $\Delta x\,\Delta y$. وعلى كل خلية، $\inf_{R_{ij}} f \cdot \Delta x \Delta y \leq
\int_{x_{i-1}}^{x_i}\int_{y_{j-1}}^{y_j} f(x,y)\,\dd y\,\dd x \leq
\sup_{R_{ij}} f \cdot \Delta x \Delta y$
برتابة التكامل بمتغير واحد (مطبقة مرتين). وبالجمع على الخلايا، يكون
التكامل المكرر $\int_a^b F$ محصورا بين مجموعي ريمان الأدنى والأعلى
للشبكة؛ وبالاتصال المنتظم يتقارب المجموعان نحو القيمة المشتركة
التي تعرّف $\iint_R f$ عندما $n \to
\infty$. وتنطبق الحجة نفسها بعد تبادل دوري
$x$ و$y$، ومنه فإن التكاملين المكررين يساويان $\iint_R f$.
\end{proof}

\begin{remark}
الاتصال على مستطيل متراص يؤدي عملا حقيقيا في برهان فوبيني: فهو
يوفر الاتصال المنتظم الذي يحصر مجاميع ريمان. أما في حالة دوال تحت
التكامل أكثر جموحا فالنص يسقط فعلا --- فهناك دوال يوجد تكاملاها
المكرران ويختلفان. والمبرهنة العامة الصادقة، مع القابلية للمكاملة
فرضية وحيدة، هي مبرهنة فوبيني من أجل تكامل لوبيغ، المبرهنة في مجلد
السنة الثالثة؛ وكل ما في هذا الفصل يبقى داخل الإطار \omterm{def:b2:metric:continuity}{المتصل} حيث
يكون البرهان الأولي أعلاه كاملا.
\end{remark}

\begin{definition}[المناطق الأولية]\label{def:b2:multint:domain}
تكون المنطقة $D \subseteq \R^2$ \emph{أولية بالنسبة إلى $y$} إذا كانت
\[
D = \{(x, y) : a \leq x \leq b,\
\varphi_1(x) \leq y \leq \varphi_2(x)\}
\]
مع $\varphi_1 \leq \varphi_2$ متصلتين على $[a,b]$ (والأولية بالنسبة إلى $x$:
بالتناظر). ومن أجل $f$ متصلة على منطقة $D$ أولية بالنسبة إلى
$y$،
\[
\iint_D f
= \int_a^b\Bigl(
\int_{\varphi_1(x)}^{\varphi_2(x)} f(x,y)\,\dd y\Bigr)\dd x ,
\]
ونتحقق (بتمديد $f$ بحجة تقريب، أو بالتجزئة) من أن التكاملين
المكررين يتطابقان عندما تكون $D$ أولية في الاتجاهين. وتعالَج
المناطق المقسومة إلى عدد منته من القطع الأولية بالجمعية، وتكون
\emph{مساحة}\index{مساحة!منطقة مستوية} $D$ هي $\operatorname
{Area}(D) = \iint_D 1$.
\end{definition}

\begin{example}\label{ex:b2:multint:triangle}
على المثلث $D = \{0 \leq x \leq 1,\ 0 \leq y \leq x\}$:
\[
\iint_D xy \,\dd x\,\dd y
= \int_0^1 x\Bigl(\int_0^x y\,\dd y\Bigr)\dd x
= \int_0^1 x\cdot\frac{x^2}{2}\,\dd x = \frac18 .
\]
وبتبديل الترتيب ($x$ من $y$ إلى $1$): $\int_0^1
y\bigl(\int_y^1 x\,\dd x\bigr)\dd y = \int_0^1
y\,\frac{1 - y^2}{2}\,\dd y = \frac18$ --- القيمة
نفسها، بحساب مختلف: فحسن اختيار ترتيب المكاملة نصف الحرفة.
\end{example}

\begin{example}[عندما ينجح ترتيب واحد فقط]\label{ex:b2:multint:swaporder}
احسب $I = \displaystyle\int_0^1\!\!\int_x^1
\eu^{y^2}\,\dd y\,\dd x$. فكما هو مكتوب، ليس للتكامل الداخلي $\int\eu^{y^2}\dd y$ دالة أصلية
أولية: ويتوقف الحساب. لكن المنطقة هي المثلث $0 \leq
x \leq y \leq 1$، وهو أولي في
الاتجاهين؛ وبتبديل الترتيب،
\[
I = \int_0^1\!\!\int_0^y \eu^{y^2}\,\dd x\,\dd y
= \int_0^1 y\,\eu^{y^2}\,\dd y
= \Bigl[\tfrac12\eu^{y^2}\Bigr]_0^1 = \frac{\eu - 1}{2} .
\]
لم يظهر المتغير الداخلي $x$ في أي موضع من الدالة تحت التكامل،
ومنه فإن مكاملته أولا أنتجت بالضبط العامل $y$ الذي يجعل التكامل
الخارجي فوريا. والعبرة: فوبيني ليست مجرد ترخيص بالتكرار --- بل هي
ترخيص \emph{بالاختيار}، والترتيب الصحيح قد يحول تكاملا مستحيلا إلى
سطر واحد. فارسم المنطقة دائما واقرأ وصفيها قبل البدء.
\end{example}

\begin{theorem}[تغيير المتغيرات]\label{thm:b2:multint:changevar}
ليكن $\Phi \colon U' \to U$ تقابلا تفاضليا من الصنف $\mathcal{C}^1$ بين مفتوحين من $\R^2$،
ولتكن $K \subseteq U$ منطقة متراصة مقسومة إلى قطع أولية مع $K' = \Phi^{-1}(K)$، ولتكن
$f$ متصلة على $K$. عندئذ
\[
\iint_K f(x, y)\,\dd x\,\dd y
= \iint_{K'} f\bigl(\Phi(u, v)\bigr)\,
\abs{\det J_\Phi(u, v)}\,\dd u\,\dd v .
\]
\end{theorem}
\admitted

\begin{remark}
البرهان الكامل --- بتقريب $\Phi$ بتفاضله على شبكة دقيقة والتحكم في
خلايا الحد --- طويل وإن لم يكن عميقا؛ وهو يُنجز كاملا في نظرية
القياس في السنة الثالثة، نتيجةً لنظرية لوبيغ. أما الحدس فهو الصورة
المستعملة أصلا من أجل \omterm{def:b2:multint:domain}{مساحة} السطوح: فالمربع الصغير ذو الضلع $\dd u$
في $(u, v)$ يُرسل، من الرتبة الأولى، على متوازي الأضلاع المولَّد
بالمتجهتين $\Phi_u\,\dd u$ و$\Phi_v\,\dd v$، ومساحته
$\abs{\det J_\Phi}\,\dd u\,\dd v$
(\cref{lem:b2:surfaces:lagrange}).
\end{remark}

\begin{remark}[طريقة: اختيار تغيير المتغيرات]
ثلاثة ردود فعل تغطي معظم الحالات. \emph{تناظر الدالة تحت التكامل}:
فالمقدار $x^2 + y^2$ يستدعي الإحداثيات القطبية، والبنية الجدائية تستدعي
الإبقاء على المحاور الديكارتية. \emph{شكل الحد}: فالحدود $u(x,y) = c_1$
و$v(x,y) = c_2$ تتوسل إلى الإحداثيات $(u, v)$ نفسها، كما في مثال المنطقة
الزائدية أدناه --- فتصير المنطقة مستطيلا، وهذا هو النصر كله.
\emph{البنية الخطية}: فالعبارات بدلالة $x + y$ و$x - y$ تدعو إلى
الدوران بالزاوية $45$ أو إلى قصّ (\cref{ex:b2:multint:affine}). وفي جميع الحالات، ثلاث
خانات ينبغي التأشير عليها قبل المكاملة: أن يكون التطبيق تقابلا من
المنطقة الجديدة على القديمة؛ وأن يُحسب جاكوبيه \emph{في الاتجاه
المستعمل فعلا} (مع القلب في النهاية إن كان أسهل)؛ وأن يدخل الجاكوبي
بقيمته المطلقة.
\end{remark}

\begin{example}[الإحداثيات القطبية]\label{ex:b2:multint:polar}
للتطبيق $\Phi(\rho, \alpha) = (\rho\cos\alpha,\ \rho\sin\alpha)$
\[
J_\Phi = \begin{pmatrix}
\cos\alpha & -\rho\sin\alpha\\
\sin\alpha & \rho\cos\alpha
\end{pmatrix},
\qquad \det J_\Phi = \rho ,
\]
ومنه $\dd x\,\dd y = \rho\,\dd\rho\,\dd\alpha$. ومن أجل القرص $D_R$ ذي نصف القطر $R$:
\[
\iint_{D_R} e^{-(x^2 + y^2)}\,\dd x\,\dd y
= \int_0^{2\pi}\!\!\int_0^R e^{-\rho^2}\rho\,\dd\rho\,\dd\alpha
= \pi\bigl(1 - e^{-R^2}\bigr)
\xrightarrow[R\to\infty]{} \pi .
\]
والمقارنة مع المربع $[-R, R]^2$ (وهو محصور بين القرصين $D_R$ و$D_{R\sqrt2}$،
وجميع الدوال تحت التكامل موجبة) تعطي $\bigl(\int_{-\infty}^\infty e^{-x^2}\dd x\bigr)^2 = \pi$:
\[
\boxed{\ \int_{-\infty}^{+\infty} e^{-x^2}\,\dd x = \sqrt{\pi}\ }
\]
--- أي تكامل غاوس من جديد، وبأشهر براهينه الآن (قارن مع الاشتقاق
ذي المتغير الواحد في \cref{ch:b2:integration}).
\end{example}

\begin{example}[تغييرات المتغيرات الأفينية]\label{ex:b2:multint:affine}
من أجل \omterm{def:b2:affine:subspace}{تطبيق أفيني} $\Phi(u, v) = M(u, v)^{\mathsf T} + C$ مع $M$ قابلة للقلب، يكون الجاكوبي هو
المصفوفة الثابتة $M$: ومنه تُضرب المساحات في العامل الثابت
$\abs{\det M}$ --- فالوعد المقطوع في \cref{ch:b2:affine} صار الآن مبرهنة.
واستعمالان فوريان. فالأهليلج $\frac{x^2}{a^2} +
\frac{y^2}{b^2} \leq 1$ هو صورة قرص الوحدة بالتطبيق
$(u, v) \mapsto (au, bv)$، ومنه فمساحته $ab \cdot \pi$ --- دون أي حساب. ومن أجل تكامل $f(x + y)$
على المربع $K = \intcc01^2$، يحول القصّ $\Phi(u, v) = (u - v, v)$ (ومحدده $1$) التكامل إلى
تكامل $f(u)$ على متوازي أضلاع، تشرّحه مبرهنة فوبيني عند $u$
ثابت: فمع $f =
\exp$،
\[
\iint_K \eu^{x+y}\,\dd x\,\dd y
= \Bigl(\int_0^1 \eu^x\,\dd x\Bigr)^2 = (\eu - 1)^2,
\]
كما تؤكد البنية الجدائية. واختيار إحداثيات ملائمة للدالة تحت
التكامل --- لا للمنطقة --- هو النصف الآخر من الحرفة.
\end{example}


\begin{example}[إحداثيات ملائمة لمنطقة منحنية الأضلاع]\label{ex:b2:multint:hyperbolicregion}
لتكن $D$ منطقة الربع الأول المحدودة بالقطعين الزائدين $xy = 1$
و$xy = 3$ وبالمستقيمين $y = x$ و$y = 3x$. وبالإحداثيات $u = xy$، $v = y/x$
تصير المنطقة المربع $\intcc13 \times \intcc13$؛ وبالقلب،
\[
x = \sqrt{u/v}, \qquad y = \sqrt{uv},
\qquad
\det J = x_uy_v - x_vy_u = \frac{1}{2v}
\]
(بحساب من سطرين مع $x = u^{1/2}v^{-1/2}$، $y =
u^{1/2}v^{1/2}$). ومنه
\[
\operatorname{Area}(D)
= \int_1^3\!\!\int_1^3\frac{\dd u\,\dd v}{2v}
= 2\cdot\frac{\ln 3}{2} = \ln 3 \approx 1.10 .
\]
أما محاولة تشريح $D$ بالإحداثيات الديكارتية فتعني تقطيعها إلى
ثلاث قطع بحدود زائدية وخطية --- وهو ممكن، لكنه بلا بهجة وكثير
الأخطاء. والعبرة تكرر \cref{ex:b2:multint:affine} بكامل قوتها: اقرأ معادلات الحد، ودع
\emph{هي} تختار الإحداثيات؛ ويحوّل الجاكوبي عندئذ \omterm{def:b2:multint:domain}{مساحة} خلية الشبكة
المنحنية الأضلاع، تماما كما فعل $\rho$ من أجل الإحداثيات القطبية.
\end{example}

\begin{example}[القيم المتوسطة]\label{ex:b2:multint:average}
\emph{القيمة المتوسطة} للدالة $f$ على منطقة $D$ هي $\frac1{\operatorname{Area}(D)}\iint_Df$.
ومثالا: متوسط المسافة عن المركز من أجل نقطة مختارة بانتظام من القرص
ذي نصف القطر $R$ هو
\[
\frac{1}{\pi R^2}\int_0^{2\pi}\!\!\int_0^R
\rho\cdot\rho\,\dd\rho\,\dd\alpha
= \frac{2\pi R^3/3}{\pi R^2} = \frac{2R}3 ,
\]
لا $R/2$: \omterm{def:b2:multint:domain}{فالمساحة} المنتظمة تضع كتلة أكبر عند أنصاف الأقطار
الكبيرة (فالتاج عند نصف القطر $\rho$ وزنه يتناسب مع $\rho$)، ومنه
يقع المتوسط بعد علامة المنتصف. وضبط هذا العامل هو بالضبط الجاكوبي
القطبي في العمل، والترجيح نفسه يفسر \omterm{def:b2:affine:barycenter}{مركز الثقل} $\bar z = 3R/8$ لنصف الكرة
المصمتة المحسوب لاحقا في هذا الفصل بدل $R/2$.
\end{example}

\section{مبرهنة غرين--ريمان}

\begin{theorem}[غرين--ريمان]\label{thm:b2:multint:green}
لتكن $K \subseteq \R^2$ منطقة متراصة أولية في الاتجاهين (أو اتحادا منتهيا لمثل
هذه المناطق ملصوقة على طول قطع)، وحدها $\partial K$ منحن مغلق من الصنف
$\mathcal{C}^1$ بالقطع وموجه \emph{عكس عقارب الساعة} (بحيث تبقى المنطقة على
اليسار). من أجل $P, Q$ من الصنف $\mathcal{C}^1$ على جوار للمنطقة $K$:
\[
\oint_{\partial K} P\,\dd x + Q\,\dd y
= \iint_K \Bigl(\frac{\partial Q}{\partial x}
- \frac{\partial P}{\partial y}\Bigr)\,\dd x\,\dd y .
\]
\end{theorem}

\begin{proof}
أولا، الطرفان جمعيان عند قطع $K$ على طول قطعة إلى جزأين $K_1, K_2$:
فالتكاملات المزدوجة تُجمع بجمعية $\iint$؛ وأما تكاملات الحد، فإن
الحدين الموجهين عكس عقارب الساعة للمنطقتين $K_1$ و$K_2$ يقطع كل
منهما القطع الداخلي مرة واحدة، في اتجاهين \emph{متعاكسين}، ومنه ففي
المجموع
\[
\oint_{\partial K_1} + \oint_{\partial K_2}
= \oint_{\partial K} + (\text{القطع، في الاتجاهين})
= \oint_{\partial K},
\]
يتلاشى المروران على طول القطع (\cref{prop:b2:multint:lineinv}) ولا يبقى إلا الحد الخارجي.
وبتكرار عدد منته من القطوع، يكفي أن نعالج منطقة أولية. ونبرهن على
$\oint P\,\dd x = -\iint_K P_y$ على منطقة $D = \{a \leq x \leq b,\ \varphi_1(x) \leq y
\leq \varphi_2(x)\}$ أولية بالنسبة إلى $y$؛ والمتطابقة
$\oint Q\,\dd y = \iint_K Q_x$ متناظرة معها (بالأولية بالنسبة إلى $x$)، والمبرهنة
مجموعهما.

نحسب التكامل المزدوج بمبرهنة فوبيني والمبرهنة الأساسية للتحليل
بمتغير واحد:
\[
\iint_D \frac{\partial P}{\partial y}\,\dd x\,\dd y
= \int_a^b \bigl(P(x, \varphi_2(x)) - P(x, \varphi_1(x))\bigr)
\,\dd x .
\]
والآن يتكون حد المنطقة $D$، عكس عقارب الساعة، مما يلي: التمثيل
البياني السفلي $y = \varphi_1(x)$ مقطوعا من اليسار إلى اليمين، والقطعة الشاقولية
اليمنى $x = b$ (صعودا)، والتمثيل البياني العلوي $y =
\varphi_2(x)$ مقطوعا
\emph{من اليمين إلى اليسار}، والقطعة الشاقولية اليسرى $x = a$
(نزولا). وعلى القطعتين الشاقوليتين يكون $x$ ثابتا، ومنه فهما
تساهمان بالمقدار $0$ في $\oint P\,\dd x$؛ أما التمثيلان البيانيان، الموسّمان
بالمتغير $x$، فيعطيان
\[
\oint_{\partial D} P\,\dd x
= \int_a^b P(x, \varphi_1(x))\,\dd x
- \int_a^b P(x, \varphi_2(x))\,\dd x
= -\iint_D \frac{\partial P}{\partial y}\,\dd x\,\dd y .
\qedhere
\]
\end{proof}

\begin{corollary}[المساحة بالحد]\label{cor:b2:multint:area}
تحت فرضيات \cref{thm:b2:multint:green}،
\[
\operatorname{Area}(K)
= \oint_{\partial K} x\,\dd y
= -\oint_{\partial K} y\,\dd x
= \frac12\oint_{\partial K} x\,\dd y - y\,\dd x .
\]
\end{corollary}

\begin{proof}
نطبق مبرهنة غرين--ريمان على $(P, Q) = (0, x)$ و$(-y, 0)$ و$\frac12(-y, x)$: ففي كل مرة
$Q_x - P_y = 1$.
\end{proof}

\begin{remark}[الاختيار بين صيغ المساحة الثلاث]
الصيغ الحدية الثلاث متساوية، لكنها ليست قابلة للتبادل عمليا. استعمل
$\oint x\,\dd y$ عندما يجعل التوسيم المقدار $\dd y$ بسيطا (كالتمثيلات البيانية
فوق المحور $y$)، واستعمل $-\oint y\,\dd x$ بالتناظر، واستعمل نصف المجموع
المتناظر عندما يعامل التوسيم $x$ و$y$ معاملة متساوية --- ففي
حالة الأهليلج أنتج دالة ثابتة تحت التكامل، دون أي تخطيط مثلثي البتة.
وعلى الحدود المضلعة يصير نصف المجموع صيغة رباط الحذاء في \cref{exo:b2:multint:12}، وهي
خوارزمية المساحين. وعندما يقطع التوسيم المعطى الحدَّ مع عقارب الساعة،
تعيد الصيغ الثلاث كلها \emph{ناقص} \omterm{def:b2:multint:domain}{المساحة}: فالنتيجة السالبة ليست
خطأ حسابيا بل تقرير توجيه --- فاقلب الإشارة، أو اقلب التوسيم.
\end{remark}

\begin{example}[مساحة الأهليلج]\label{ex:b2:multint:ellipsearea}
من أجل $x = a\cos t$، $y = b\sin t$، $t \in [0, 2\pi]$:
\[
\operatorname{Area}
= \frac12\int_0^{2\pi}\bigl(a\cos t \cdot b\cos t
- b\sin t\cdot(-a\sin t)\bigr)\,\dd t
= \frac{ab}{2}\int_0^{2\pi}\dd t = \pi ab .
\]
\end{example}

\begin{example}[غرين--ريمان تحققا متقاطعا]\label{ex:b2:multint:greencheck}
خذ $P = -y^3$، $Q = x^3$ على قرص الوحدة المغلق $D$. من جهة الحد، مع
$\gamma(t) = (\cos t, \sin t)$:
\[
\oint_{\partial D}P\,\dd x + Q\,\dd y
= \int_0^{2\pi}\bigl(\sin^4 t + \cos^4 t\bigr)\dd t
= 2\pi\cdot\Bigl(\frac38 + \frac38\Bigr) = \frac{3\pi}2 ,
\]
بالتخطيط ($\sin^4 + \cos^4 = \tfrac34 +
\tfrac14\cos4t$). ومن جهة الداخل:
\[
\iint_D(Q_x - P_y)\,\dd x\,\dd y
= \iint_D 3(x^2 + y^2)\,\dd x\,\dd y
= 3\int_0^{2\pi}\!\!\int_0^1\rho^3\,\dd\rho\,\dd\alpha
= \frac{3\pi}2 .
\]
العدد نفسه، بحسابين مختلفين جدا --- وهذا هو الاستعمال العملي: فأيهما
كان أسهل من طرفي متطابقة غرين صار هو الحساب، والآخر تحققا. وفي حالة
دورانات الحقول كثيرة الحدود حول منحنيات مغلقة، يكون التكامل المزدوج
هو الطرف الأسهل دائما تقريبا.
\end{example}

\begin{remark}
تفسّر مبرهنة غرين--ريمان \cref{ex:b2:multint:angleform}: فمن أجل \omterm{def:b2:multint:exact}{صورة مغلقة} ($Q_x = P_y$)،
ينعدم التكامل حول حد أي منطقة محتواة في $U$. ولا تفشل صورة
الزاوية في أن تكون تامة إلا لأن الثقب في المبدأ يمنع القرص المحدود
بدائرة الوحدة من الوقوع داخل $U$ --- فالتكاملات المنحنية للصور
المغلقة تكشف ثقوب المنطقة. (وبدفع هذه الملاحظة أبعد، تصير كوهومولوجيا
دي رام.)
\end{remark}

\section{التكاملات الثلاثية}

تمتد النظرية إلى ثلاثة متغيرات دون أي فكرة جديدة: فمبرهنة فوبيني
تختزل $\iiint$ إلى ثلاثة تكاملات ذات متغير واحد (إما \emph{بالتشريح}:
$\iiint_K f = \int\bigl(\iint_{K_z}
f\bigr)\dd z$ على الشرائح الأفقية $K_z$، وإما \emph{بالتكديس}: بالمكاملة
بدلالة $z$ أولا على طول أعمدة شاقولية)، وتتحقق صيغة تغيير
المتغيرات مع الجاكوبي من الحجم $3 \times 3$.

\begin{example}[الإحداثيات الأسطوانية والكروية]\label{ex:b2:multint:spherical}
\emph{الأسطوانية} $(x, y, z) = (\rho\cos\alpha, \rho\sin\alpha,
z)$: $\dd x\,\dd y\,\dd z = \rho\,\dd\rho\,\dd\alpha\,\dd z$.
\emph{الكروية} $(x, y, z) = (r\cos\theta\cos\varphi,\
r\sin\theta\cos\varphi,\ r\sin\varphi)$ ($\theta$ الطول، $\varphi \in [-\frac\pi2, \frac\pi2]$ العرض): بنشر \omterm{def:b2:linalg:det}{المحدد}
من الحجم $3
\times 3$ وفق السطر الأخير،
\[
\det J = r^2\cos\varphi ,
\qquad
\dd x\,\dd y\,\dd z
= r^2\cos\varphi\;\dd r\,\dd\theta\,\dd\varphi .
\]
\omterm{pb:b2:multint:1}{وحجم الكرة المصمتة} ذات نصف القطر $R$:
\[
V = \int_0^R\!\!\int_0^{2\pi}\!\!\int_{-\pi/2}^{\pi/2}
r^2\cos\varphi\;\dd\varphi\,\dd\theta\,\dd r
= \frac{R^3}{3}\cdot 2\pi \cdot 2
= \boxed{\frac43\pi R^3} ,
\]
وبذلك تُسدَّد أخيرا الصيغة المسلَّم بها في فصول الحجوم من الكتب
السابقة.
\end{example}

\begin{example}[رباعي الوجوه، مرتين]
\label{ex:b2:multint:tetra}
حجم $T = \{x, y, z \geq 0,\ x + y + z \leq 1\}$، بالتكديس: من أجل $(x, y)$ ثابت في المثلث $x + y \leq 1$، يجري
$z$ على $\intcc0{1 - x - y}$، ومنه
\[
V = \int_0^1\!\!\int_0^{1-x}(1 - x - y)\,\dd y\,\dd x
= \int_0^1\frac{(1 - x)^2}{2}\,\dd x = \frac16 .
\]
وبالتشريح: المقطع عند الارتفاع $z$ هو المثلث $\{x, y
\geq 0,\ x + y \leq 1 - z\}$، ومساحته
$\frac{(1-z)^2}2$، ومنه $V = \int_0^1\frac{(1-z)^2}2\,\dd z = \frac16$ من جديد --- فالحسابان هما التكاملان نفساهما
بترتيب مختلف، وهذا كل ما تدّعيه مبرهنة فوبيني. والقيمة $\frac16 =
\frac13\cdot\frac12\cdot1$ هي
صيغة المخروط (\cref{ex:b2:multint:cone}) بقاعدة مثلثية، والصيغة ذات البعد $n$
وهي $1/n!$ يُبرهن عليها بهذا التشريح بالضبط في مسألة نهاية الأسبوع.
\end{example}

\begin{example}[مركز ثقل نصف الكرة المصمتة]\label{ex:b2:multint:centroid}
من أجل نصف الكرة المصمتة العلوي $H$ ذي نصف القطر $R$ ($z \geq 0$)،
يكون ارتفاع \omterm{def:b2:affine:barycenter}{مركز الثقل} $\bar z = \frac1{V}\iiint_H z$، مع $V =
\frac23\pi R^3$. وبالإحداثيات الكروية
($z =
r\sin\varphi$، $\varphi \in \intcc0{\pi/2}$):
\[
\iiint_H z
= \int_0^R r^3\,\dd r\int_0^{2\pi}\dd\theta
\int_0^{\pi/2}\sin\varphi\cos\varphi\,\dd\varphi
= \frac{R^4}4\cdot2\pi\cdot\frac12 = \frac{\pi R^4}4 ,
\]
ومنه
\[
\bar z = \frac{\pi R^4/4}{2\pi R^3/3} = \frac{3R}8 :
\]
فنقطة توازن نصف الكرة المصمتة تقع على ثلاثة أثمان نصف القطر فوق
الوجه المستوي --- أي تحت نصف الارتفاع $R/2$، كما يجب، لأن الجسم
أسمن قرب القاعدة. ولكل حساب مركز ثقل هذا الشكل: تكامل عزم واحد،
وحجم واحد، ونسبة واحدة، وتحقق من المعقولية بمقارنتها بالهندسة.
\end{example}

\begin{example}[الحجم بالتشريح: المخروط]\label{ex:b2:multint:cone}
مخروط \omterm{def:b2:multint:domain}{مساحة} قاعدته $A$ وارتفاعه $h$ (رأسه إلى الأعلى وقاعدته
عند $z = 0$): المقطع عند الارتفاع $z$ هو القاعدة مضروبة في
المعامل $(1 -
z/h)$، ومساحته $A(1 - z/h)^2$. ومنه
\[
V = \int_0^h A\Bigl(1 - \frac zh\Bigr)^2\dd z = \frac{Ah}{3} :
\]
أي الثلث في صيغ المدرسة، وهي صالحة من أجل \emph{أي} شكل قاعدة ---
فالتشريح يحولها إلى تكامل مربع.
\end{example}

\begin{example}[عتبات القابلية للمكاملة في المستوي]\label{ex:b2:multint:threshold}
من أجل أي $\alpha > 0$ يتقارب $\iint_{D}\rho^{-\alpha}\,\dd
x\,\dd y$ على قرص الوحدة المثقوب $D$
(بالنهاية على التيجان $\varepsilon \leq \rho \leq 1$)؟ بالإحداثيات القطبية،
\[
\int_0^{2\pi}\!\!\int_\varepsilon^1\rho^{-\alpha}\,
\rho\,\dd\rho\,\dd\alpha
= 2\pi\int_\varepsilon^1\rho^{1-\alpha}\,\dd\rho ,
\]
وهو يتقارب عندما $\varepsilon \to 0$ إذا وفقط إذا كان $1 - \alpha > -1$، أي $\alpha < 2$: ففي
البعد $2$ يكون أسّ الشذوذ الحرج هو البعد نفسه، والعامل الإضافي
$\rho$ الآتي من الجاكوبي يلطّف الشذوذ بقوة واحدة. (وبالمثل $\alpha < 3$
من أجل شذوذ نقطي في الفضاء، عبر $r^2$.) ومسك الحسابات القطري من
هذا النوع هو كيف يتقرر أمر القابلية للمكاملة في لمحة في إطار لوبيغ
في السنة الثالثة --- وهو سبب تقارب $\iiint 1/r$ بلا عناء في \cref{exo:b2:multint:7}.
\end{example}

\begin{remark}[مزالق شائعة]
(1) \emph{التوجيه}: يتغير \omterm{def:b2:multint:lineint}{التكامل المنحني} إشارةً بتغير اتجاه السير،
وتقتضي مبرهنة غرين--ريمان أن يكون الحد عكس عقارب الساعة (بالمنطقة
على اليسار)؛ وفي حالة منطقة ذات ثقب، يُقطع الحد الداخلي \emph{مع
عقارب الساعة}. (2) \emph{يدخل الجاكوبي بقيمته المطلقة}: فتغيير
المتغيرات لا ينتج أبدا \omterm{def:b2:multint:domain}{مساحة} سالبة، ونسيان $\abs{\det}$ يقلب الإشارات
عادة بالضبط عندما يعكس التطبيق التوجيه. (3) \emph{العامل القطبي
$\rho$}: $\dd x\,\dd y = \rho\,\dd\rho\,\dd\alpha$، لا $\dd\rho\,\dd\alpha$ --- وهو أكثر الأخطاء شيوعا في الفصل
كله؛ ويمسكه التحليل البعدي، لأن $\dd\rho\,\dd\alpha$ بُعده بُعد طول لا بُعد
\omterm{def:b2:multint:domain}{مساحة}. (4) \emph{التكاملات المزدوجة المعتلة}: تتوافق النهايات على
الأقراص المتنامية وعلى المربعات المتنامية هنا لأن الدوال تحت التكامل
موجبة (بالحصر)؛ أما في حالة دوال تتغير إشارتها فقد تتعلق النهاية
بالاستنفاد المختار، ولا يُدّعى شيء دون تقارب مطلق. (5) \emph{المناطق
مقابل الدوال تحت التكامل}: دالة جدائية تحت التكامل على منطقة غير
جدائية \emph{لا} تفكك التكامل --- فالتفكيك يحتاج إلى الاثنين معا،
كما في مربع \cref{ex:b2:multint:polar}.
\end{remark}

\begin{remark}[آفاق داخل هذا المجلد]
تكامل غاوس المحسوب هنا موجود بصمت في كل مكان من فصول الاحتمالات:
فالثابت $\sqrt\pi$ داخل صيغة ستيرلنغ (\cref{thm:b2:comparison:stirling}) هو تكامل هذا الفصل، وهو
يحدد عبر ستيرلنغ السلوك المقارب $1/\sqrt{\pi n}$ لاحتمالات عودة السير
العشوائي في \cref{ch:b2:proba}. وتعود \omterm{pb:b2:multint:1}{تكاملات واليس} في مسألة نهاية الأسبوع
هناك أيضا، فتقود تقديرات المعامل الثنائي المركزي نفسها. وفي الاتجاه
الآخر، يكمل عنصرا \omterm{def:b2:multint:domain}{المساحة} والحجم في هذا الفصل هندسة \cref{ch:b2:surfaces}، وتعيد
صيغة غرين حساب مساحات المغلِّفات في \cref{ch:b2:curves} (\omterm{pb:b2:curves:1}{النجمية}، في \cref{exo:b2:multint:5}).
ففصل واحد، وثلاث خدمات: قياس من أجل الهندسة، وثوابت من أجل
الاحتمالات، وانضباط تغيير المتغيرات الذي يستعمله الاثنان.
\end{remark}

\section{تمارين}

\begin{exercise}[$\star$]\label{exo:b2:multint:1}
احسب $\int_\gamma y^2\,\dd x + x\,\dd y$ على: (1) القطعة من $(0,0)$ إلى $(1,1)$؛ (2) قوس القطع
المكافئ $y = x^2$ من $(0,0)$ إلى $(1,1)$. وهل \omterm{def:b2:multint:exact}{الصورة تامة}؟
\end{exercise}

\begin{exercise}[$\star$]\label{exo:b2:multint:2}
بيّن أن $\omega = (2xy + y^3)\,\dd x + (x^2 + 3xy^2 + 1)\,\dd y$ مغلقة على $\R^2$، وجد كمونا، واحسب $\int_\gamma\omega$ على أي
قوس من $(0, 0)$ إلى $(1, 2)$.
\end{exercise}

\begin{exercise}[$\star$]\label{exo:b2:multint:3}
احسب $\iint_D (x + y)\,\dd x\,\dd y$ حيث $D$ هي المنطقة المحدودة بالمنحنيين $y = x^2$
و$y = x$ ($0 \leq x \leq 1$)، بترتيبي المكاملة كليهما.
\end{exercise}

\begin{exercise}[$\star\star$]\label{exo:b2:multint:4}
باستعمال الإحداثيات القطبية، احسب $\iint_D \frac{\dd x\,\dd y}{(1 +
x^2 + y^2)^2}$ على المستوي كله (بوصفه
نهاية على الأقراص)، واحسب $\iint_{D'} xy\,\dd x\,\dd y$ على ربع القرص $D' = \{x, y
\geq 0,\ x^2 + y^2 \leq 1\}$.
\end{exercise}

\begin{exercise}[$\star\star$]\label{exo:b2:multint:5}
احسب \omterm{def:b2:multint:domain}{المساحة} التي تحصرها \omterm{pb:b2:curves:1}{النجمية} $x = \cos^3 t$، $y =
\sin^3 t$، $t \in [0, 2\pi]$، باستعمال
\cref{cor:b2:multint:area}. \emph{(خطّط $\sin^2 t\cos^2 t$.)}
\end{exercise}

\begin{exercise}[$\star\star$]\label{exo:b2:multint:6}
احسب حجم المجسم المحدود من الأسفل بالمجسم المكافئ $z
= x^2 + y^2$ ومن الأعلى
بالمستوي $z = 1$، بالطريقتين: بالتكديس (بمكاملة $1 - x^2 - y^2$ على قرص
الوحدة، بالإحداثيات القطبية) وبالتشريح (فالشرائح الأفقية أقراص نصف
قطرها $\sqrt z$).
\end{exercise}

\begin{exercise}[$\star\star$]\label{exo:b2:multint:7}
(جذب كرة مصمتة --- مبرهنة نيوتن، حالة خاصة) بيّن أن حجم القشرة
الكروية $a \leq r \leq
b$ هو $\frac43\pi(b^3 - a^3)$، واحسب $\iiint_{B}
\frac{\dd x\,\dd y\,\dd z}{r}$ على الكرة المصمتة $B$
ذات نصف القطر $R$ ($r$ المسافة إلى المبدأ). \emph{(بالإحداثيات
الكروية.)}
\end{exercise}

\begin{exercise}[$\star\star\star$]\label{exo:b2:multint:8}
(مبرهنة بوانكاريه، الحالة \omterm{pb:b2:curves:1}{النجمية}) ليكن $U$ نجمي الشكل بالنسبة
إلى $0$ (أي $M \in U \Rightarrow [0, M] \subseteq U$) ولتكن $\omega = P\dd x + Q\dd y$ \omterm{def:b2:multint:exact}{صورة مغلقة} من الصنف $\mathcal{C}^1$
على $U$. نعرّف
\[
f(x, y) = \int_0^1 \bigl(x\,P(tx, ty) + y\,Q(tx, ty)\bigr)\dd t .
\]
باستعمال المفاضلة تحت علامة التكامل (\cref{ch:b2:integration}) و$P_y = Q_x$، بيّن أن
$f_x = P$ و$f_y = Q$: فكل \omterm{def:b2:multint:exact}{صورة مغلقة} على مفتوح نجمي الشكل تكون تامة.
\end{exercise}

\begin{exercise}[$\star\star\star$]\label{exo:b2:multint:9}
(تكامل ديريكليه بالمكاملة المزدوجة) برر واستغل
\[
\int_0^\infty\!\!\int_0^\infty e^{-xy}\sin x\;\dd y\,\dd x
\quad\text{مقابل}\quad
\int_0^\infty\!\!\int_0^\infty e^{-xy}\sin x\;\dd x\,\dd y
\]
على $[0, A] \times [0, \infty)$: بيّن $\int_0^A \frac{\sin
x}{x}\dd x = \frac\pi2 - \int_0^\infty
e^{-Ay}\frac{y\sin A + \cos A}{1 + y^2}\dd y$ واستعد $\int_0^\infty \frac{\sin x}{x}\,\dd x = \frac\pi2$، بالمقارنة مع برهان التكامل
ذي الوسيط في \cref{ch:b2:integration}.
\end{exercise}

\begin{exercise}[$\star\star\star$]\label{exo:b2:multint:10}
(متراجحة تساوي المحيط عبر فيرتنغر) ليكن $\gamma$ منحنى مغلقا بسيطا
من الصنف $\mathcal{C}^1$ طوله $2\pi$، موسّما \omterm{def:b2:curves:length}{بطول القوس} على $[0, 2\pi]$، ويحصر
\omterm{def:b2:multint:domain}{مساحة} $A$. باستعمال \cref{cor:b2:multint:area} ومبرهنة بارسيفال ومتراجحة فيرتنغر
(من تمارين \cref{ch:b2:fourier})، برهن على $A \leq \pi$، مع التساوي من أجل الدائرة.
\emph{(وحّد $\int_0^{2\pi} x(s)\dd s
= 0$؛ واكتب $2A = \oint x\,\dd y - y\,\dd x$ واحصر $2A \leq
\int (x^2 + y'^2)$ بعناية عبر $2A = \int_0^{2\pi}(xy' - yx')\dd
s$
و$x^2 + y'^2 \geq 2xy'$.)}
\end{exercise}

\begin{exercise}[$\star\star$]\label{exo:b2:multint:11}
(عزوم الكرة المصمتة) من أجل الكرة المصمتة $B$ ذات نصف القطر
$R$ في $\R^3$، احسب $\iiint_B z^2\,\dd x\,\dd y\,\dd z$ بالإحداثيات الكروية، واستنتج $\iiint_B (x^2 + y^2 +
z^2)\,\dd x\,\dd y\,\dd z$
بالتناظر. وتحقق متقاطعا من الأخير مع حساب القشرة $\int_0^R r^2\cdot4\pi
r^2\,\dd r$.
\end{exercise}

\begin{exercise}[$\star\star$]\label{exo:b2:multint:12}
(صيغة رباط الحذاء) ليكن $K$ مضلعا رؤوسه $(x_1,
y_1), \dots, (x_m, y_m)$ مرتبة عكس عقارب
الساعة (بالأدلة بترديد $m$). استنتج من \cref{cor:b2:multint:area} أن
\[
\operatorname{Area}(K)
= \frac12\sum_{i=1}^m
\bigl(x_iy_{i+1} - x_{i+1}y_i\bigr) ,
\]
وتحقق من الصيغة على المثلث $(0,0)$، $(1,0)$، $(0,1)$.
\end{exercise}

\section{مسألة: حجم الكرة المصمتة في البعد $n$}

\begin{problem}[{مسألة نهاية الأسبوع --- $V_n =
\pi^{n/2}/\Gamma(\frac n2 + 1)$، وغرابة الأبعاد
العليا}]\label{pb:b2:multint:1}
\omterm{def:b2:multint:domain}{مساحة} القرص $\pi$، \omterm{pb:b2:multint:1}{وحجم الكرة المصمتة} $\frac43\pi$ --- ثم ماذا؟ تحسب
هذه المسألة حجم كرة الوحدة المصمتة في $\R^n$ من أجل كل $n$،
مرتين (بتراجع تشريحي تقوده \omterm{pb:b2:multint:1}{تكاملات واليس}، ثم عبر الدالة $\Gamma$
وتكامل غاوس في \cref{ex:b2:multint:polar})، ثم تقرأ الهندسة: فالحجوم تبلغ ذروتها في
البعد الخامس وتندفع إلى الصفر، ويختبئ كل حجم الكرة العالية البعد
تقريبا في قشرة رقيقة قرب حدها.\index{حجم الكرة المصمتة}\index{تكاملات واليس}
ومن أجل دالة متصلة على كرة مصمتة من $\R^n$، يُفهم التكامل تكاملا
مكررا $n$ مرة (بتشريح إحداثية في كل مرة، كما في الفصل من أجل
$n
\leq 3$)؛ ونكتب $B_n(R)$ للكرة المصمتة المغلقة ذات نصف القطر $R$
ومركزها $0$، و$v_n(R)$ لحجمها، و$V_n =
v_n(1)$، مع $V_0 = 1$ بالاصطلاح.

\medskip
\noindent\textbf{الجزء الأول --- التراجع التشريحي.}
\begin{enumerate}
  \item بالتعويض $x_i = Ru_i$ في كل من التكاملات المكررة $n$، بيّن
        $v_n(R) = V_nR^n$.
  \item بتشريح $B_n(1)$ على طول إحداثيتها الأخيرة، بيّن
        \[
        V_n = V_{n-1}\int_{-1}^{1}(1 -
        t^2)^{\frac{n-1}2}\,\dd t .
        \]
  \item مع $t = \sin\theta$، حدد التكامل بوصفه تكامل واليس: $\int_{-1}^1(1 -
        t^2)^{\frac{n-1}2}\dd t = 2W_n$، حيث
        $W_n =
        \int_0^{\pi/2}\cos^n\theta\,\dd\theta =
        \int_0^{\pi/2}\sin^n\theta\,\dd\theta$.
  \item برهن على متطابقتي واليس (بالمكاملة بالتجزئة؛ ثم بالتلسكوب
        $nW_nW_{n-1}$):
        \[
        W_n = \frac{n-1}nW_{n-2} \quad (n \geq 2),
        \qquad
        W_nW_{n-1} = \frac{\pi}{2n} \quad (n \geq 1).
        \]
\end{enumerate}

\medskip
\noindent\textbf{الجزء الثاني --- حل التراجع.}
\begin{enumerate}[resume]
  \item اجمع الأسئلة من 2 إلى 4 في التراجع ذي الخطوتين
        \[
        V_n = \frac{2\pi}{n}\,V_{n-2}
        \qquad (n \geq 2).
        \]
  \item استنتج الصيغ المغلقة، من أجل $k \geq 0$:
        \[
        V_{2k} = \frac{\pi^k}{k!},
        \qquad
        V_{2k+1} = \frac{2^{k+1}\pi^k}{1\cdot3\cdot5\cdots
        (2k+1)} .
        \]
  \item جدول $V_1, \dots, V_7$ عدديا. وباستعمال النسبة $V_n/V_{n-2} = 2\pi/n$ وقيمتي $2W_5$
        و$2W_6$، برهن على أن المتتالية $(V_n)$ متزايدة حتى قيمتها
        العظمى $V_5 = \frac{8\pi^2}{15} \approx
        5.26$ ومتناقصة من ثم فصاعدا.
  \item بيّن أن $V_n \to 0$ أسرع من أي متتالية هندسية، وأن $\sum_{n\geq1} V_n$
        متقارب: فمجموع حجوم كرات الوحدة المصمتة كلها منته.
  \item برهن على المتطابقة المولّدة
        \[
        \sum_{k\geq0} V_{2k}\,x^{2k} = \eu^{\pi x^2}
        \qquad (x \in \R),
        \]
        واستنتج $\sum_{k\geq0}V_{2k} = \eu^\pi \approx
        23.14$.
\end{enumerate}

\medskip
\noindent\textbf{الجزء الثالث --- الطريق الثاني: $\Gamma$ وتكامل
غاوس.}
\begin{enumerate}[resume]
  \item بيّن بمبرهنة فوبيني (فالدالة تحت التكامل جدائية) أن
        \[
        I_n = \int_{\R^n}\eu^{-\norm
        x^2}\dd x
        = \Bigl(\int_{-\infty}^{+\infty}
        \eu^{-t^2}\dd t\Bigr)^{\!n} = \pi^{n/2},
        \]
        وهو تكامل غاوس ذو البعد $n$، مفهوما نهايةً على المكعبات
        $\intcc{-R}{R}^n$.
  \item ذكّر بأن $\Gamma(s) = \int_0^\infty t^{s-1}\eu^{-t}\dd
        t$ (\cref{def:b2:integration:gamma}). وانطلاقا من $\Gamma(s+1) = s\,\Gamma(s)$ (\cref{thm:b2:integration:gammaprops})
        و$\Gamma(\tfrac12) = \sqrt\pi$ (بالتعويض $t = u^2$ واستدعاء تكامل غاوس)، احسب
        \[
        \Gamma(k + 1) = k!,
        \qquad
        \Gamma\Bigl(k + \frac32\Bigr) =
        \frac{1\cdot3\cdots(2k+1)}{2^{k+1}}\,\sqrt\pi .
        \]
  \item برهن، بالتراجع عبر تراجع السؤال 5، على الصيغة الوحيدة
        \[
        V_n = \frac{\pi^{n/2}}{\Gamma\bigl(\frac n2 +
        1\bigr)} \qquad (n \geq 1),
        \]
        وتحقق من أنها تعيد إنتاج الصيغتين المغلقتين في السؤال 6.
  \item بيّن $\int_0^\infty \eu^{-r^2}r^{n-1}\dd r =
        \tfrac12\Gamma\bigl(\tfrac n2\bigr)$ واستنتج المتطابقة
        \[
        I_n = n\,V_n\int_0^\infty \eu^{-r^2}\,r^{n-1}\,\dd r.
        \]
        وفسّرها: فالكتلة الغاوسية للفضاء $\R^n$ تُجمع على طول قشور
        كروية تكون ``مساحتها ذات البعد $(n-1)$'' عند نصف القطر
        $r$ هي $nV_nr^{n-1}$ --- وقد بُرهن على الطرفين الآن برهانا
        مستقلا، ومنه فالتفسير لا يكلف شيئا.
  \item ضع $s_{n-1} = nV_n$ (وهي \omterm{def:b2:multint:domain}{مساحة} كرة الوحدة $S^{n-1}$، بانسجام مع $v_n(R) = \int_0^R
        s_{n-1}r^{n-1}\dd r$).
        جدول $s_0, \dots, s_3$ وتحقق من $s_1 = 2\pi$ و$s_2 = 4\pi$ و$s_3 = 2\pi^2$.
\end{enumerate}

\medskip
\noindent\textbf{الجزء الرابع --- الأبعاد العليا غريبة.}
\begin{enumerate}[resume]
  \item انطلاقا من صيغة ستيرلنغ (\cref{thm:b2:comparison:stirling}) مطبقة على $k!$، بيّن
        من أجل $n = 2k$ زوجي:
        \[
        V_n \sim \frac{1}{\sqrt{\pi n}}
        \Bigl(\frac{2\pi\eu}{n}\Bigr)^{n/2}
        \qquad (n \to \infty, \ n \text{ زوجي}),
        \]
        واشرح لماذا يمتد حصر التلاشي فوق الهندسي نفسه إلى $n$
        الفردي عبر التراجع.
  \item تقع كرة الوحدة المصمتة داخل المكعب $\intcc{-1}1^n$ ذي الحجم $2^n$.
        احسب نسبة الملء $V_n/2^n$ من أجل $n = 2, 3, 10$، وبيّن أنها تؤول إلى
        $0$: ففي البعد العالي، يقع كل المكعب تقريبا في زواياه.
  \item بيّن أن نسبة $v_n(1)$ الواقعة على مسافة لا تتجاوز $\varepsilon$ من
        كرة الحد هي $1 -
        (1 - \varepsilon)^n \to 1$؛ وعدديا، أي نسبة من كرة مصمتة ذات بعد
        $100$ تقع في القشرة الخارجية ذات السماكة $1\%$؟
  \item برهن على السلوك المقارب لواليس $W_n \sim
        \sqrt{\dfrac{\pi}{2n}}$ \emph{(برتابة $(W_n)$،
        والنسبة $W_n/W_{n-2} \to 1$، و$W_nW_{n-1}
        = \frac\pi{2n}$)}، وعلى الحصر السفلي $W_n \geq
        \sqrt{\dfrac{\pi}{2(n+1)}}$ من أجل
        كل $n$.
  \item (التركيز على شريحة) نسبة كرة الوحدة المصمتة التي إحداثيتها
        الأولى تتجاوز $\delta$ هي $\int_\delta^1(1 - x^2)^{\frac{n-1}2}\dd x \,\big/\,
        (2W_n)$. باستعمال $1 - u \leq \eu^{-u}$ والحصر
        الذيلي $\int_\delta^\infty \eu^{-a x^2}\dd x \leq
        \frac{\eu^{-a\delta^2}}{2a\delta}$، بيّن أن هذه النسبة لا تتجاوز
        \[
        \frac{\eu^{-(n-1)\delta^2/2}}{(n-1)\,\delta}
        \Big/ \sqrt{\frac{2\pi}{n+1}}
        \]
        واستنتج: من أجل $\delta = s/\sqrt{n-1}$، تقع الكرة المصمتة كلها عدا نسبة
        $O(\eu^{-s^2/2}/s)$ منها في الشريحة $\abs{x_1} \leq s/\sqrt{n-1}$. فالكرة المصمتة العالية
        البعد، إحصائيا، فطيرة رقيقة في كل اتجاه في آن واحد.
  \item اجمع الأسئلة من 16 إلى 19 في فقرة واحدة: أين يقع حجم $B_n(1)$
        (قرب كرة الحد، ومع ذلك داخل شرائح $O(1/\sqrt n)$ لكل مستو زائد يمر
        بالمركز)، ولماذا لا يتناقض النصان.
\end{enumerate}

\medskip
\noindent\textbf{الجزء الخامس --- أجسام أخرى، وتركيب.}
\begin{enumerate}[resume]
  \item (المبسّط) لتكن $\Delta_n = \{x \in \R^n : x_i \geq 0,\
        \sum x_i \leq 1\}$. بيّن بالتشريح وبالتراجع أن $\operatorname{vol}(\Delta_n) = \frac1{n!}$.
  \item (متعدد الوجوه المتقاطع) استنتج أن حجم $C_n = \{x :
        \sum\abs{x_i} \leq 1\}$ هو $\frac{2^n}{n!}$،
        وتحقق من الحصر $C_n \subseteq B_n(1)
        \subseteq \intcc{-1}1^n$ على مستوى الحجوم: $\frac{2^n}{n!} \leq V_n \leq 2^n$.
  \item احسب $V_4$ بطريقة ثالثة: شرّح $\R^4 = \R^2 \times
        \R^2$، وكامل \omterm{def:b2:multint:domain}{مساحة} القرص
        $(z, w)$ على القرص $(x, y)$ بالإحداثيات القطبية، واستعد $V_4 = \frac{\pi^2}2$.
  \item (مونت كارلو في ورطة) تُسحب نقطة بانتظام في المكعب $\intcc{-1}1^{20}$.
        بيّن أن احتمال وقوعها في الكرة المصمتة المرسومة داخله هو
        $V_{20}/2^{20} \approx 2.5\cdot10^{-8}$، ومنه فإن نحو أربعين مليون سحبة لازمة قبل أن يُنتظر
        أول إصابة: فتقدير $V_n$ بالمعاينة بالرفض ينهار في البعد
        العالي (وهي لعنة الأبعاد).
  \item تركيب. التقى اشتقاقان مستقلان عند $V_n =
        \pi^{n/2}/\Gamma(\frac n2 + 1)$: اسرد أي مبرهنة
        من هذا الفصل استعمل كل منهما (فوبيني، تغيير المتغيرات،
        تكامل غاوس القطبي)، وأي معطيات ذات متغير واحد (واليس،
        $\Gamma$، ستيرلنغ). وأين يعيد مجلد السنة الثالثة هذا الحساب
        بنظرية لوبيغ، وماذا يضيف؟
\end{enumerate}
\end{problem}
